\chapter{svnserve Reference---Custom Subversion
Server}\label{svn.ref.svnserve}

\section{svnserve}\label{svn.ref.svnserve.re}

\emph{Serve Subversion repositories via Subversion\textquotesingle s
custom network protocol}

\textbf{Synopsis}

\texttt{svnserve\ {[}-d\ \textbar{}\ -i\ \textbar{}\ -t\ \textbar{}\ -X{]}\ OPTIONS...}

\subsection{Description}\label{svn.ref.svnserve.re.desc}

\texttt{svnserve} allows access to Subversion repositories using
Subversion\textquotesingle s custom network protocol.

You can run \texttt{svnserve} as a standalone server process (for
clients that are using the \texttt{svn://} access method); you can have
a daemon such as \texttt{inetd} or \texttt{xinetd} launch it for you on
demand (also for \texttt{svn://}), or you can have \texttt{sshd} launch
it on demand for the \texttt{svn+ssh://} access method.

Unless the \texttt{-\/-config-file} option was specified on the command
line, once the client has selected a repository by transmitting its URL,
\texttt{svnserve} reads a file named \texttt{conf/svnserve.conf} in the
repository directory to determine repository-specific settings such as
what authentication database to use and what authorization policies to
apply. See \hyperref[svn.serverconfig.svnserve]{svnserve, a Custom
Server} for details of the \texttt{svnserve.conf} file.

\subsection{Options}\label{svn.ref.svnserve.re.options}

Unlike the previous commands we\textquotesingle ve described,
\texttt{svnserve} has no subcommands---it is controlled exclusively by
options.

\begin{description}
\item[\protect\phantomsection\label{svn.ref.svnserve.sw.cache_fulltexts}{}\texttt{-\/-cache-fulltexts}
\textless ARG\textgreater{}]
Toggles support for fulltext file content caching (in FSFS repositories
only).
\item[\protect\phantomsection\label{svn.ref.svnserve.sw.cache_txdeltas}{}\texttt{-\/-cache-txdeltas}
\textless ARG\textgreater{}]
Toggles support for file content delta caching (in FSFS repositories
only).
\item[\protect\phantomsection\label{svn.ref.svnserve.sw.compression}{}\texttt{-\/-compression}
\textless LEVEL\textgreater{}]
Specifies the level of compression used for wire transmissions as an
integer beween 0 and 9, inclusive. A value of \texttt{9} offers the best
compression, \texttt{5} is the default value, and \texttt{0} disables
compression altogether.
\item[\protect\phantomsection\label{svn.ref.svnserve.sw.config_file}{}\texttt{-\/-config-file}
\textless FILENAME\textgreater{}]
When specified, \texttt{svnserve} reads \textless FILENAME\textgreater{}
once at program startup and caches the \texttt{svnserve} configuration.
The password and authorization configurations referenced from filename
will be loaded on each connection. \texttt{svnserve} will not read any
per‐repository \texttt{conf/svnserve.conf} files when this option is
used. See the \hyperref[svn.serverconfig.svnserve]{svnserve, a Custom
Server} for details of the file format for this option.
\item[\protect\phantomsection\label{svn.ref.svnserve.sw.daemon}{}\texttt{-\/-daemon}
(\texttt{-d})]
Causes \texttt{svnserve} to run in daemon mode. \texttt{svnserve}
backgrounds itself and accepts and serves TCP/IP connections on the
\texttt{svn} port (3690, by default).
\item[\protect\phantomsection\label{svn.ref.svnserve.sw.foreground}{}\texttt{-\/-foreground}]
When used together with \texttt{-d}, causes \texttt{svnserve} to stay in
the foreground. This is mainly useful for debugging.
\item[\protect\phantomsection\label{svn.ref.svnserve.sw.inetd}{}\texttt{-\/-inetd}
(\texttt{-i})]
Causes \texttt{svnserve} to use the \texttt{stdin} and \texttt{stdout}
file descriptors, as is appropriate for a daemon running out of
\texttt{inetd}.
\item[\protect\phantomsection\label{svn.ref.svnserve.sw.help}{}\texttt{-\/-help}
(\texttt{-h})]
Displays a usage summary and exits.
\item[\protect\phantomsection\label{svn.ref.svnserve.sw.listen_host}{}\texttt{-\/-listen-host}
\textless HOST\textgreater{}]
Causes \texttt{svnserve} to listen on the interface specified by
\textless HOST\textgreater, which may be either a hostname or an IP
address.
\item[\protect\phantomsection\label{svn.ref.svnserve.sw.listen_once}{}\texttt{-\/-listen-once}
(\texttt{-X})]
Causes \texttt{svnserve} to accept one connection on the \texttt{svn}
port, serve it, and exit. This option is mainly useful for debugging.
\item[\protect\phantomsection\label{svn.ref.svnserve.sw.listen_port}{}\texttt{-\/-listen-port}
\textless PORT\textgreater{}]
Causes \texttt{svnserve} to listen on \textless PORT\textgreater{} when
run in daemon mode. (FreeBSD daemons listen only on tcp6 by
default---this option tells them to also listen on tcp4.)
\item[\protect\phantomsection\label{svn.ref.svnserve.sw.log_file}{}\texttt{-\/-log-file}
\textless FILENAME\textgreater{}]
Instructs \texttt{svnserve} to create (if necessary) and use the file
located at \textless FILENAME\textgreater{} for Subversion operational
log output of the same sort that \texttt{mod\_dav\_svn} generates. See
\hyperref[svn.serverconfig.operational-logging]{High-level Logging} for
details.
\item[\protect\phantomsection\label{svn.ref.svnserve.sw.memory_cache_size}{}\texttt{-\/-memory-cache-size}
(\texttt{-M}) \textless ARG\textgreater{}]
Configures the size (in Megabytes) of the extra in-memory cache used to
minimize redundant operations. The default value is \texttt{16}. (This
cache is used for FSFS-backed repositories only.)
\item[\protect\phantomsection\label{svn.ref.svnserve.sw.pid_file}{}\texttt{-\/-pid-file}
\textless FILENAME\textgreater{}]
Causes \texttt{svnserve} to write its process ID to
\textless FILENAME\textgreater, which must be writable by the user under
which \texttt{svnserve} is running.
\item[\protect\phantomsection\label{svn.ref.svnserve.sw.prefer_ipv6}{}\texttt{-\/-prefer-ipv6}
(\texttt{-6})]
When resolving the listen hostname, prefer an IPv6 answer over an IPv4
one. IPv4 is preferred by default.
\item[\protect\phantomsection\label{svn.ref.svnserve.sw.quiet}{}\texttt{-\/-quiet}]
Disables progress notifications. Error output will still be printed.
\item[\protect\phantomsection\label{svn.ref.svnserve.sw.root}{}\texttt{-\/-root}
(\texttt{-r}) \textless ROOT\textgreater{}]
Sets the virtual root for repositories served by \texttt{svnserve}. The
pathname in URLs provided by the client will be interpreted relative to
this root and will not be allowed to escape this root.
\item[\protect\phantomsection\label{svn.ref.svnserve.sw.threads}{}\texttt{-\/-threads}
(\texttt{-T})]
When running in daemon mode, causes \texttt{svnserve} to spawn a thread
instead of a process for each connection (e.g., for when running on
Windows). The \texttt{svnserve} process still backgrounds itself at
startup time.
\item[\protect\phantomsection\label{svn.ref.svnserve.sw.tunnel}{}\texttt{-\/-tunnel}
(\texttt{-t})]
Causes \texttt{svnserve} to run in tunnel mode, which is just like the
\texttt{inetd} mode of operation (both modes serve one connection over
\texttt{stdin}/\texttt{stdout}, and then exit), except that the
connection is considered to be preauthenticated with the username of the
current UID. This flag is automatically passed for you by the client
when running over a tunnel agent such as \texttt{ssh}. That means
there\textquotesingle s rarely any need for \emph{you} to pass this
option to \texttt{svnserve}. So, if you find yourself typing
\texttt{svnserve\ -\/-tunnel} on the command line and wondering what to
do next, see \hyperref[svn.serverconfig.svnserve.sshauth]{Tunneling over
SSH}.
\item[\protect\phantomsection\label{svn.ref.svnserve.sw.tunnel_user}{}\texttt{-\/-tunnel-user}
\textless NAME\textgreater{}]
Used in conjunction with the \texttt{-\/-tunnel} option, tells
\texttt{svnserve} to assume that \textless NAME\textgreater{} is the
authenticated user, rather than the UID of the \texttt{svnserve}
process. This is useful for users wishing to share a single system
account over SSH, but to maintain separate commit identities.
\item[\protect\phantomsection\label{svn.ref.svnserve.sw.version}{}\texttt{-\/-version}]
Displays version information and a list of repository backend modules
available, and then exits.
\end{description}

